\section{Requisitos Funcionais}

% Garante numeracao ate o nivel de subsubsection no documento
\setcounter{secnumdepth}{3}

Os requisitos funcionais estão organizados em grupos identificados por siglas que em conjunto formam identificadores únicos. Ator e prioridade são
declarados no nível do grupo e herdados pelos requisitos; exceções são indicadas no próprio requisito. Os critérios de aceite são apresentados
como lista de condições verificáveis.

\subsection{Comportamentos Transversais dos Cadastros}

Para promover consistência e reduzir repetições, os seguintes comportamentos transversais aplicam-se a todos os requisitos de manutenção de cadastros (REQ-CAD-*), salvo indicação explícita em contrário:

\begin{itemize}[leftmargin=0.6cm]
    \item \textbf{Confirmação explícita para operações críticas:} O sistema deve exigir confirmação explícita do usuário para operações de inativação, reativação, cancelamento ou exclusão lógica, quando aplicáveis.
    \item \textbf{Mensagens de erro associadas a campos:} O sistema deve exibir mensagem de erro junto ao campo que causou a rejeição, sem apagar os demais campos do formulário.
    \item \textbf{Preservação de dados em formulários:} O sistema deve preservar os valores informados nos campos do formulário durante validações e tentativas de salvamento, exceto quando explicitamente indicado.
    \item \textbf{Atualização imediata de consultas:} Após operação concluída com sucesso, o sistema deve refletir imediatamente o resultado nas listagens e consultas relevantes.
    \item \textbf{Pré-preenchimento em edição:} O sistema deve apresentar formulários de edição pré-preenchidos com os dados existentes do registro, garantindo consistência e facilitando correções.
    \item \textbf{Impedimento de exclusão física com vínculos:} O sistema deve impedir exclusão física de registros com vínculos associados, exibindo mensagem explicativa.
    \item \textbf{Manutenção de dados após inativação:} O sistema deve manter o cadastro de um registro inativo consultável com seus dados íntegros, sem exclusão física ou ocultação do estado.
    \item \textbf{Seleção controlada de associações:} Para campos que referenciam outras entidades, o sistema deve permitir seleção apenas a partir de opções válidas e ativas apresentadas, impedindo digitação livre e associação a registros inativos em novos vínculos quando aplicável.
    \item \textbf{Imutabilidade do identificador interno:} O sistema deve manter o identificador interno do registro único e inalterável após a criação, preservando-o durante qualquer operação de atualização.
    \item \textbf{Disponibilidade em consultas:} O sistema deve recuperar registros ativos e inativos nas consultas de cadastro, salvo indicação contrária no requisito.
\end{itemize}

Estes comportamentos são herdados implicitamente pelos requisitos de cadastro, sendo omitidos das listas de critérios de aceite individuais para evitar repetição. Exceções ou especificidades são indicadas explicitamente no requisito correspondente.

% Contadores hierárquicos internos para os identificadores RF
% O requisito incorpora a sigla do macrogrupo atual e o numero
% do grupo atual, acrescentando apenas sua sequencia local.
\newcounter{rfmacrogrupo}
\newcounter{rfgrupo}[rfmacrogrupo]
\newcounter{rfreq}[rfgrupo]
\newcounter{rfcadastrotopico}[rfmacrogrupo]
\newcommand{\rfsiglasecao}{RF}
\newcommand{\rfmacrosigla}{}
% \newcommand{\rfgrupoid}{\thesection.\arabic{rfmacrogrupo}.\arabic{rfgrupo}}
\newcommand{\rfgrupoid}{\thesubsection.\arabic{rfgrupo}}
% \newcommand{\rfcadastrotopicoid}{\arabic{rfmacrogrupo}.\arabic{rfcadastrotopico}}
%% \newcommand{\rfcadastrotopicoid}{\thesubsection.\arabic{rfcadastrotopico}}
\newcommand{\rfcadastrotopicoid}{%
  \ifnum\value{rfgrupo}=0
    \thesubsection.\arabic{rfcadastrotopico}%
  \else
    \thesubsection.\arabic{rfgrupo}.\arabic{rfcadastrotopico}%
  \fi
}
\renewcommand{\therfreq}{\rfgrupoid.\arabic{rfreq}}
\newcommand{\rfreqid}{\rfsiglasecao-\rfmacrosigla-\therfreq}

\newcommand{\rfmacroprioridadebase}{}
\newcommand{\rfmacroprioridadepadrao}{}

\makeatletter
\def\reqcad@split#1 #2\relax{%
  \gdef\reqcad@code{#1}%
  \gdef\reqcad@title{#2}%
}
\gdef\reqcad@isreq{0}
\def\reqcad@testREQ#1#2#3#4\@nil{%
  \if R#1%
    \if E#2%
      \if Q#3%
        \gdef\reqcad@isreq{1}%
      \else
        \gdef\reqcad@isreq{0}%
      \fi
    \else
      \if F#2%
        \if -#3%
          \gdef\reqcad@isreq{1}%
        \else
          \gdef\reqcad@isreq{0}%
        \fi
      \else
        \gdef\reqcad@isreq{0}%
      \fi
    \fi
  \else
    \gdef\reqcad@isreq{0}%
  \fi
}
\makeatother

\newcommand{\rfmacrogrupo}[6]{%
  \refstepcounter{rfmacrogrupo}%
  \setcounter{rfgrupo}{0}%
  \gdef\rfmacrosigla{#2}%
  \gdef\rfmacroprioridadebase{#5}%
  \gdef\rfmacroprioridadepadrao{#5}%
  \subsection{#3 (#2)}%
  \label{#1}%
  \textbf{Escopo:} #6\\
  \textbf{Ator:} #4\\
  \textbf{Prioridade:} #5

  \vspace{0.3cm}
}

\newcommand{\rfgruporeq}[2]{%
  \refstepcounter{rfgrupo}%
  \setcounter{rfreq}{0}% -<<<<<<<<<<<<<<<<<<<<<<<<<<<<<
  \setcounter{rfcadastrotopico}{0}%
  \gdef\rfmacroprioridadepadrao{\rfmacroprioridadebase}%
  \subsubsection{#2 (\rfmacrosigla)}%
  \label{#1}%
  \vspace{0.3cm}
}

\newcommand{\rfgruporeqator}[3]{%
  \refstepcounter{rfgrupo}%
  \setcounter{rfreq}{0}%
  \gdef\rfmacroprioridadepadrao{\rfmacroprioridadebase}%
  \subsubsection{#2 (\rfmacrosigla)}%
  \label{#1}%
  \noindent\textbf{Ator:} #3

  \vspace{0.3cm}
}

\newcommand{\rfgruporeqatorprio}[4]{%
  \refstepcounter{rfgrupo}%
  \setcounter{rfreq}{0}%
  \gdef\rfmacroprioridadepadrao{#4}%
  \subsubsection{#2 (\rfmacrosigla)}%
  \label{#1}%
  \noindent\textbf{Ator:} #3\\
  \textbf{Prioridade:} #4

  \vspace{0.3cm}
}

\newcommand{\rfgruposemrequisitos}[1]{%
  \noindent\textit{#1}

  \vspace{0.3cm}
}

\newcommand{\rfcadastrotopico}[3]{%
  \refstepcounter{rfcadastrotopico}%
  \gdef\rfmacroprioridadepadrao{#3}%
  \noindent\textbf{\rfcadastrotopicoid\hspace{0.5em}#1}\par
  \noindent\textbf{Código base:} #2\\
  \textbf{Prioridade:} #3

  \vspace{0.2cm}
}

\newcommand{\rfmostraprioridade}[1]{%
  \begingroup
  \edef\rfprioridaderequisito{#1}%
  \edef\rfprioridademacro{\rfmacroprioridadepadrao}%
  \ifx\rfprioridaderequisito\rfprioridademacro
  \else
    \textbf{Prioridade específica:} #1\\
  \fi
  \endgroup
}

% #1 = rotulo de referencia (vazio para requisitos REQ-*)
% #2 = nome do requisito funcional
% #3 = descricao
% #4 = prioridade
% #5 = dependencia
% #6 = criterio de aceite
\makeatletter
\newcommand{\reqfuncfull}[6]{%
  \refstepcounter{rfreq}%
  \reqcad@split#2 \relax
  \expandafter\reqcad@testREQ\reqcad@code ZZZ\@nil
  \ifnum\reqcad@isreq=1\relax
    \edef\@currentlabel{\reqcad@code}%
    \noindent\textbf{[\reqcad@code] \reqcad@title}\\
    \label{\reqcad@code}%
    \textbf{Descrição:} #3\\
    \textbf{Dependência:} #5\\
    \textbf{Critério de aceite:} #6
  \else
    \noindent\textbf{[\rfreqid] #2}\\
    \label{#1}%
    \textbf{Descrição:} #3\\
    \rfmostraprioridade{#4}%
    \textbf{Dependência:} #5\\
    \textbf{Critério de aceite:} #6
  \fi

  \vspace{0.4cm}
}
\makeatother

\newcommand{\reqfunc}[5]{%
  \reqfuncfull{#1}{#2}{#3}{\rfmacroprioridadepadrao}{#4}{#5}%
}

\newcommand{\reqfuncprio}[6]{%
  \reqfuncfull{#1}{#2}{#3}{#4}{#5}{#6}%
}

\newcommand{\reqfunccad}[5]{%
  \reqfuncfull{}{#1}{#2}{#3}{#4}{#5}%
}

\rfmacrogrupo
{mg:autenticacao_perfil}
{AP}
{Autenticação e Perfil}
{Definido no nível do grupo}
{Essencial}
{Organizar a autenticação, a recuperação de acesso e a identificação do perfil de acesso no sistema.}
\rfgruporeqator
{g:acesso_identidade_aluno}
{Acesso e identidade do aluno}
{Aluno}

\reqfunc
{rf:autenticar_login_aluno}
{Autenticar login do aluno}
{O sistema deve autenticar o acesso do aluno.}
{Depende do requisito [\ref{REQ-CAD-ALU-01}].}
{\begin{itemize}[leftmargin=0.6cm]
    \item O sistema deve aceitar CPF ou e-mail como identificador de login do aluno.
    \item O sistema deve autenticar o aluno antes de liberar funcionalidades protegidas, impedindo acesso com credenciais inválidas, ausentes, expiradas ou não vinculadas a perfil de aluno.
    \item O sistema deve informar o motivo da falha de autenticação de forma compatível com o cenário ocorrido, sem revelar dados de outros perfis.
\end{itemize}}

\reqfunc
{rf:recuperar_senha_aluno}
{Recuperar senha do aluno}
{O sistema deve redefinir a credencial de acesso do aluno mediante verificação de identidade.}
{Depende do requisito [\ref{REQ-CAD-ALU-01}].}
{\begin{itemize}[leftmargin=0.6cm]
    \item O sistema deve aceitar e-mail como identificador para recuperação de senha do aluno.
    \item O sistema deve permitir iniciar o fluxo sem autenticação.
    \item Ao confirmar a identidade, o sistema deve exigir o CPF e a definição de nova senha antes do acesso.
    \item O sistema deve recusar identificações inválidas, inexistentes ou bloqueadas, e rejeitar validações incorretas, expiradas ou já utilizadas, sem revelar dados.
\end{itemize}}

\reqfunc
{rf:recuperar_acesso_aluno_dispositivo}
{Recuperar acesso do aluno em novo dispositivo}
{O sistema deve recuperar o acesso do aluno à conta em novo dispositivo, mediante verificação de identidade.}
{Depende do requisito [\ref{REQ-CAD-ALU-01}].}
{\begin{itemize}[leftmargin=0.6cm]
    \item O sistema deve permitir recuperação mediante confirmação de identidade.
    \item Após recuperação, o sistema deve invalidar sessões anteriores.
    \item O sistema deve impedir acesso e recusar validações inválidas, expiradas ou inconsistentes enquanto a identidade não for confirmada.
\end{itemize}}

\reqfunc
{rf:restringir_acesso_perfil_aluno}
{Restringir acesso do perfil de aluno}
{O sistema deve restringir o acesso do aluno às funcionalidades autorizadas.}
{Depende do requisito [\ref{rf:autenticar_login_aluno}].}
{\begin{itemize}[leftmargin=0.6cm]
    \item O sistema deve permitir acesso apenas às funcionalidades autorizadas ao aluno.
    \item O sistema deve negar acesso a funcionalidades fora do escopo do perfil.
    \item O sistema deve exibir mensagem de permissão insuficiente quando necessário.
    \item O sistema não deve expor dados administrativos ou de outros perfis.
\end{itemize}}

\reqfunc
{rf:direcionar_login_aluno}
{Direcionar login do aluno}
{O sistema deve direcionar o aluno autenticado diretamente ao seu dashboard, sem opção de seleção de perfil.}
{Depende do requisito [RF-AP-\ref{rf:autenticar_login_aluno}].}
{\begin{itemize}[leftmargin=0.6cm]
    \item O sistema deve redirecionar automaticamente para o dashboard do aluno após autenticação bem-sucedida.
    \item O sistema não deve apresentar opções de seleção de perfil para alunos.
\end{itemize}}
\rfgruporeqator
{g:acesso_identidade_professora}
{Acesso e identidade da professora}
{Professora}

\reqfunc
{rf:autenticar_login_professora}
{Autenticar login da professora}
{O sistema deve autenticar o acesso da professora.}
{Nenhuma.}
{\begin{itemize}[leftmargin=0.6cm]
    \item O sistema deve aceitar CPF ou e-mail como identificador de login da professora.
    \item O sistema deve autenticar a professora antes de liberar funcionalidades protegidas, impedindo acesso com credenciais inválidas, ausentes, expiradas ou não vinculadas a perfil de professora.
    \item O sistema deve exibir mensagem de rejeição compatível com o cenário de falha, sem expor dados protegidos do sistema.
\end{itemize}}

\reqfunc
{rf:restringir_acesso_perfil_professora}
{Restringir acesso do perfil de professora}
{O sistema deve restringir o acesso da professora às funcionalidades autorizadas.}
{Depende do requisito [RF-AP-\ref{rf:autenticar_login_professora}].}
{\begin{itemize}[leftmargin=0.6cm]
    \item O sistema deve restringir o perfil de professora às funcionalidades autorizadas: gestão administrativa e de repertório, manutenção de vínculos entre turma e \textit{mix}, visualização operacional nominal da turma, cancelamento administrativo de \textit{check-ins} e consulta dos painéis gerenciais.
    \item O sistema deve negar acesso a funcionalidades não declaradas para a professora, exibindo mensagem de permissão insuficiente ou de contexto não autorizado.
    \item O sistema não deve disponibilizar ações além das previstas para o perfil de professora.
\end{itemize}}

\reqfunc
{rf:direcionar_login_professora}
{Direcionar login da professora}
{O sistema deve apresentar à professora opções de acesso após autenticação, permitindo seleção entre participar como aluno, gestão de aula ou gestão administrativa.}
{Depende do requisito [RF-AP-\ref{rf:autenticar_login_professora}].}
{\begin{itemize}[leftmargin=0.6cm]
    \item O sistema deve oferecer três opções: participar como aluno (\textit{dashboard} do aluno), gestão de aula (\textit{dashboard} da professora) e gestão administrativa (\textit{dashboard} administrativo).
    \item O sistema deve permitir seleção da opção desejada pela professora.
    \item Após seleção, o sistema deve direcionar para o \textit{dashboard} correspondente.
\end{itemize}}

\rfmacrogrupo
{mg:checkin_aluno}
{CA}
{Check-in do Aluno}
{Aluno}
{Essencial}
{Organizar a jornada do aluno na consulta, reserva e gestão do próprio \textit{check-in}.}

\rfgruporeqator
{g:jornada_checkin_aluno}
{Jornada do aluno no check-in}
{Aluno}

\reqfunc
{rf:painel_aulas_dia}
{Painel geral de aulas do dia}
{O sistema deve exibir ao aluno um painel com as aulas disponíveis no dia atual.}
{Depende dos requisitos [\ref{REQ-CAD-TUR-01}] e [\ref{REQ-CAD-SAL-01}].}
{\begin{itemize}[leftmargin=0.6cm]
    \item O sistema deve disponibilizar apenas turmas da data atual.
    \item O sistema deve fornecer, para cada turma disponível, horário de início, horário de fim, professora, sala, total de vagas reserváveis e vagas ocupadas.
    \item O sistema deve desconsiderar turmas inativas, sem horário definido, sem sala válida associada ou sem professora associada.
    \item O sistema deve excluir a \textit{bike} da professora do total de vagas reserváveis.
    \item O sistema deve apresentar mensagem informativa quando não houver turmas disponíveis no dia.
\end{itemize}}

\reqfunc
{rf:detalhe_aula_selecionada}
{Exibir detalhes da aula selecionada}
{O sistema deve apresentar ao aluno os detalhes da aula selecionada.}
{Depende dos requisitos [RF-CA-\ref{rf:painel_aulas_dia}], [\ref{REQ-CAD-BIK-01}] e [\ref{REQ-CAD-MAN-01}].}
{\begin{itemize}[leftmargin=0.6cm]
    \item O sistema deve exibir horário de início, horário de fim, nome da professora e sala da turma selecionada.
    \item O sistema deve exibir os dados da turma selecionada sem misturar informações de outra turma.
    \item O sistema deve indicar a presença ou ausência de \textit{mix} vigente.
    \item Quando houver \textit{mix} vigente, o sistema deve disponibilizar acesso ao repertório e aos detalhes musicais associados à aula.
    \item O sistema deve disponibilizar o total de vagas reserváveis, vagas disponíveis e a representação das posições associadas às \textit{bikes}.
    \item O sistema deve classificar cada posição conforme seu estado operacional (`livre', `ocupada', `em manutenção', `reservada para a professora'), representar visualmente cada estado de forma distinta, e permitir seleção apenas de posições `livre'.
    \item O sistema deve identificar a posição associada ao aluno autenticado quando houver \textit{check-in} ativo.
    \item O sistema deve permitir ao aluno visualizar apenas o primeiro nome dos alunos com \textit{check-in} ativo nas posições ocupadas, sem expor sobrenome, dados de contato ou outras informações de identificação.
    \item O sistema deve excluir \textit{bikes} em manutenção e a posição da professora das vagas reserváveis e disponíveis.
    \item Se a turma deixar de estar disponível, ficar sem sala válida ou perder dados mínimos antes da abertura do painel, o sistema deve informar a indisponibilidade e retornar ao painel geral sem exibir informações desatualizadas.
\end{itemize}}

\reqfunc
{rf:realizar_checkin}
{Realizar check-in}
{O sistema deve registrar o check-in do aluno autenticado em uma vaga reservável da turma selecionada, desde que atendidas as condições de elegibilidade, disponibilidade e consistência da operação.}
{Depende dos requisitos [\ref{REQ-CAD-ALU-01}] e [RF-CA-\ref{rf:detalhe_aula_selecionada}].}
{\begin{itemize}[leftmargin=0.6cm]
    \item O sistema deve permitir o registro do \textit{check-in} somente dentro da janela de 30 minutos antes do \texttt{horarioInicio} até o próprio \texttt{horarioInicio} da turma.
    \item O sistema deve permitir reserva apenas em posição classificada como livre e reservável, recusando posições indisponíveis, incluindo aquelas associadas à professora ou em manutenção.
    \item O sistema deve impedir a realização de mais de um \textit{check-in} ativo para o mesmo aluno na mesma turma e data, bem como em turmas com sobreposição de horário.
    \item O sistema deve permitir o registro apenas quando a sessão autenticada corresponder ao aluno para o qual a reserva está sendo realizada.
    \item O sistema deve recusar a operação quando a vaga, a turma ou a janela de reserva deixarem de estar disponíveis durante o fluxo.
    \item Após registro bem-sucedido, o sistema deve atualizar imediatamente a representação das posições, a contagem de vagas e indicar a confirmação da operação ao aluno.
\end{itemize}}

\reqfunc
{rf:cancelar_checkin}
{Cancelar check-in}
{O sistema deve cancelar o check-in ativo do aluno quando solicitado por ele.}
{Depende do requisito [RF-CA-\ref{rf:realizar_checkin}].}
{\begin{itemize}[leftmargin=0.6cm]
    \item O sistema deve aceitar cancelamento apenas do \textit{check-in} ativo do próprio aluno autenticado e dentro do prazo até o \texttt{horarioInicio} da turma, recusando cancelamentos de reservas de terceiros ou fora do prazo com mensagem apropriada.
    \item O sistema deve liberar a vaga imediatamente após o cancelamento bem-sucedido.
    \item O sistema deve manter o \textit{check-in} cancelado visível com status de cancelado.
    \item Quando a reserva já não estiver ativa no momento da confirmação, o sistema deve informar o estado atualizado sem exibir mensagem de sucesso.
\end{itemize}}

\reqfunc
{rf:fila_espera}
{Entrar na fila de espera}
{O sistema deve registrar o aluno na fila de espera da turma quando ela estiver lotada. A fila serve como registro de interesse do aluno; a gestão de quem efetivamente ocupa uma vaga liberada é feita presencialmente pela professora, fora do sistema.}
{Depende dos requisitos [\ref{REQ-CAD-ALU-01}], [RF-CA-\ref{rf:detalhe_aula_selecionada}] e [RF-CA-\ref{rf:realizar_checkin}].}
{\begin{itemize}[leftmargin=0.6cm]
    \item O sistema deve aceitar no máximo 5 alunos na fila da turma, recusando novas entradas quando cheia e exibindo mensagem informativa.
    \item Se a turma deixar de estar lotada antes da conclusão da entrada, não deve registrar o aluno na fila.
    \item O sistema deve impedir duplicidade de entrada na fila para a mesma turma e data.
    \item O sistema deve impedir coexistência de fila de espera e \textit{check-in} confirmado para a mesma turma e data.
    \item O sistema deve impedir entrada da professora na fila.
    \item O sistema deve exibir a fila ordenada por posição de entrada no mapa operacional da professora, para apoio à gestão presencial.
\end{itemize}}

\reqfuncprio
{rf:consultar_repertorio_aula}
{Consultar repertório e detalhes musicais da aula}
{O sistema deve exibir ao aluno o repertório do \textit{mix} vigente e os detalhes musicais associados à aula.}
{Importante}
{Depende dos requisitos [\ref{REQ-CAD-MIX-01}], [\ref{REQ-CAD-TMX-01}] e [RF-CA-\ref{rf:detalhe_aula_selecionada}].}
{\begin{itemize}[leftmargin=0.6cm]
    \item O sistema deve disponibilizar as músicas do \textit{mix} vigente na sequência planejada.
    \item O sistema deve permitir acesso aos detalhes de cada música do repertório.
    \item Ao detalhar uma música, o sistema deve exibir os dados disponíveis de artista ou banda e categoria musical.
    \item Ao detalhar uma música, o sistema deve exibir os links de vídeo-aula associados quando houver associação válida com \textit{VideoAula}.
    \item O sistema deve sinalizar ausência de artista, categoria ou link apenas no item afetado, sem impedir a consulta das demais músicas.
    \item Quando não houver músicas no \textit{mix}, o sistema deve informar a ausência de repertório disponível.
    \item Se o \textit{mix} deixar de estar disponível ao abrir a lista, o sistema deve informar a indisponibilidade.
\end{itemize}}

\reqfuncprio
{rf:acesso_painel_participacao}
{Acesso ao painel de participação do aluno}
{O sistema deve disponibilizar ao aluno, no fluxo principal, um acesso ao painel de participação do aluno.}
{Desejável}
{Depende dos requisitos [RF-AP-\ref{rf:autenticar_login_aluno}] e [RF-AA-\ref{rf:consultar_historico_de_presenca}].}
{\begin{itemize}[leftmargin=0.6cm]
    \item O sistema deve permitir ao aluno autenticado acesso ao próprio painel de participação, encaminhando-o somente ao painel vinculado à sessão autenticada.
    \item O ponto de acesso deve exibir apenas um indicativo de entrada ao painel de acompanhamento individual.
    \item Quando não houver dados de participação, o painel de destino deve abrir com mensagem informativa.
\end{itemize}}

\reqfuncprio
{rf:avaliar_musica_repertorio}
{Registrar e modificar avaliação de música}
{O sistema deve permitir que o aluno autenticado registre ou modifique sua avaliação para uma música disponível no repertório.}
{Desejável}
{Depende dos requisitos [RF-CA-\ref{rf:consultar_repertorio_aula}] e [\ref{REQ-CAD-MUS-01}].}
{\begin{itemize}[leftmargin=0.6cm]
    \item O sistema deve permitir que o aluno registre uma avaliação para cada música disponível no repertório consultado.
    \item O sistema deve permitir a alteração posterior da avaliação atribuída pelo próprio aluno.
    \item O sistema deve manter apenas uma avaliação vigente por combinação de aluno e música.
    \item Ao consultar uma música previamente avaliada, o sistema deve apresentar a avaliação atualmente registrada pelo aluno.
    \item O sistema deve impedir que o aluno registre ou modifique avaliações em nome de outro usuário.
    \item O sistema deve informar o resultado da operação após o registro ou a atualização da avaliação.
\end{itemize}}



\rfmacrogrupo
{mg:gestao_administrativa}
{GADM}
{Gestão Administrativa}
{Professora}
{Essencial}
{Manter a base cadastral e a estrutura operacional que sustentam a operação das aulas.}

Nos subtópicos padronizados de cadastro, o item `Padronização de entradas' deve ser declarado apenas quando houver regra específica de enumeração, seleção a partir de listas controladas, cardinalidade de associação ou controle formal de entrada relevante para a entidade. Quando não houver restrição dessa natureza, o item pode ser omitido nos novos cadastros; ajustes de rastreabilidade histórica devem ser registrados em documento auxiliar, sem aparecer no documento de entrega.

\rfcadastrotopico{Manter Aluno}{REQ-CAD-ALU}{Essencial}

\reqfunccad
{REQ-CAD-ALU-01 Manter aluno}
{O sistema deve manter o cadastro de alunos, que são os participantes das turmas de spinning e cujo estado ativo determina a elegibilidade para check-in.}
{Essencial}
{Nenhuma.}
{\begin{itemize}[leftmargin=0.6cm]
    \item O sistema deve permitir cadastrar, consultar, atualizar, inativar e reativar alunos.
    \item O sistema deve exigir `nome', `cpf', `email', `dataNascimento' e `genero' como campos obrigatórios.
    \item O sistema deve permitir `telefone', `urlFoto', `instagram', `facebook', `tiktok' e `observacoes' como campos opcionais.
    \item O sistema deve validar formato de `cpf', `email' e `urlFoto' quando informada, unicidade de `cpf' e `email', consistência de `dataNascimento', e limitar `genero' às opções \texttt{MASCULINO}, \texttt{FEMININO} e \texttt{OUTRO} na inclusão e atualização.
    \item O sistema deve permitir consulta por `nome', `cpf', `email' e status, com ordenação alfabética por nome.
    \item O sistema deve preservar o histórico de participação do aluno após inativação.
    \item O sistema deve recusar a inativação quando houver tratamento operacional obrigatório pendente.
\end{itemize}}

\rfcadastrotopico{Manter Grupo de Alunos}{REQ-CAD-GAL}{Importante}

\reqfunccad
{REQ-CAD-GAL-01 Manter grupo de alunos}
{O sistema deve manter o cadastro de grupos de alunos, que organizam alunos em conjuntos para fins de gestão e comunicação operacional.}
{Importante}
{Depende do requisito [\ref{REQ-CAD-ALU-01}].}
{\begin{itemize}[leftmargin=0.6cm]
    \item O sistema deve permitir cadastrar, consultar, atualizar, inativar e reativar grupos de alunos.
    \item O sistema deve exigir `nome' como campo obrigatório nos fluxos de criação e edição.
    \item O sistema deve permitir criar grupo sem alunos associados inicialmente.
    \item O sistema deve validar unicidade de `nome' na inclusão e na atualização.
    \item O sistema deve permitir consulta por `nome' e status, com ordenação alfabética por nome.
    \item O sistema deve exibir a composição atual do grupo quando houver alunos associados.
    \item O sistema deve impedir que o mesmo aluno seja associado mais de uma vez ao mesmo grupo.
    \item O sistema deve manter o cadastro individual dos alunos inalterado ao incluir ou remover alunos do grupo.
    \item O sistema deve manter inalterado o estado cadastral dos alunos associados quando o grupo for inativado.
\end{itemize}}

\rfcadastrotopico{Manter Fabricante}{REQ-CAD-FAB}{Importante}

\reqfunccad
{REQ-CAD-FAB-01 Manter fabricante}
{O sistema deve manter o cadastro de fabricantes, que identificam a origem das bikes e são referenciados no cadastro de equipamentos.}
{Importante}
{Nenhuma.}
{\begin{itemize}[leftmargin=0.6cm]
    \item O sistema deve permitir cadastrar, consultar, atualizar, inativar e reativar fabricantes.
    \item O sistema deve exigir `nome' como campo obrigatório.
    \item O sistema deve permitir `emailContato' e `telefoneContato' como campos opcionais.
    \item O sistema deve validar a unicidade de `nome' na inclusão e na atualização.
    \item O sistema deve permitir consulta por `nome', por dados de contato informados e por status, com ordenação alfabética por nome.
\end{itemize}}

\rfcadastrotopico{Manter Tipo de Manutenção}{REQ-CAD-TMA}{Importante}

\reqfunccad
{REQ-CAD-TMA-01 Manter tipo de manutenção}
{O sistema deve manter o cadastro de tipos de manutenção, que classificam as intervenções realizadas nas bikes e determinam o registro operacional correspondente.}
{Importante}
{Nenhuma.}
{\begin{itemize}[leftmargin=0.6cm]
    \item O sistema deve permitir cadastrar, consultar, atualizar, inativar e reativar tipos de manutenção.
    \item O sistema deve exigir `nome' como campo obrigatório.
    \item O sistema deve permitir `descricao' como campo opcional.
    \item O sistema deve validar a unicidade de `nome' na inclusão e na atualização.
    \item O sistema deve permitir consulta por `nome' e por status, com ordenação alfabética por nome.
    \item O sistema deve preservar a exibição do tipo originalmente associado em manutenções já registradas, mesmo após inativação.
\end{itemize}}

\rfcadastrotopico{Manter Sala}{REQ-CAD-SAL}{Essencial}

\reqfunccad
{REQ-CAD-SAL-01 Manter sala}
{O sistema deve manter o cadastro de salas, que definem o espaço físico e a grade de posições onde as bikes são alocadas para uso nas turmas.}
{Essencial}
{Nenhuma.}
{\begin{itemize}[leftmargin=0.6cm]
    \item O sistema deve permitir cadastrar, consultar, atualizar, inativar e reativar salas.
    \item O sistema deve exigir `nome', `numeroFilas', `numeroColunas', `numeroFilaProfessora' e `numeroColunaProfessora' como campos obrigatórios nos fluxos de inclusão e atualização.
    \item O sistema deve apresentar um layout visual das posições das bikes na sala ao definir `numeroFilas' e `numeroColunas', representando cada posição como uma bike disponível.
    \item O sistema deve aceitar apenas valores válidos para a grade da sala, exigindo `numeroFilas' e `numeroColunas' positivos, e `numeroFilaProfessora' e `numeroColunaProfessora' dentro dos limites da grade física.
    \item O sistema deve definir a posição da professora (formada por `numeroFilaProfessora' e `numeroColunaProfessora') como a bike que fica inativa para check-in de alunos, pois é utilizada pela professora durante a aula.
    \item O sistema deve impedir que a grade física dependa de listas previamente cadastradas de outras entidades.
    \item O sistema deve permitir consulta por `nome' e status, com ordenação alfabética por nome.
    \item O sistema deve recusar alterações à grade da sala que invalidem posições de bike existentes ou sejam incompatíveis com turmas ativas associadas.
\end{itemize}}

\rfcadastrotopico{Manter Bike}{REQ-CAD-BIK}{Essencial}

\reqfunccad
{REQ-CAD-BIK-01 Manter bike}
{O sistema deve manter o cadastro de bikes, que representam os equipamentos físicos de spinning associados a fabricantes e posicionados nas salas para uso nas turmas.}
{Essencial}
{Depende dos requisitos [\ref{REQ-CAD-FAB-01}] e [\ref{REQ-CAD-SAL-01}].}
{\begin{itemize}[leftmargin=0.6cm]
    \item O sistema deve permitir cadastrar, consultar, atualizar, inativar e reativar bikes.
    \item O sistema deve permitir cadastrar bike sem associação física imediata à sala.
    \item O sistema deve exigir `fabricante', `nome' e `dataCadastro' como campos obrigatórios.
    \item O sistema deve validar formato e consistência cronológica de `dataCadastro'.
    \item O sistema deve validar unicidade de `nome' na inclusão e na atualização.
    \item O sistema deve permitir consulta por `nome', `fabricante' e status, com ordenação alfabética por nome.
    \item O sistema deve permitir associar ou reposicionar bike apenas em posição existente, livre e compatível com a grade da sala.
    \item O sistema deve impedir que a mesma bike ocupe mais de uma posição simultaneamente.
\end{itemize}}

\rfcadastrotopico{Manter Turma}{REQ-CAD-TUR}{Essencial}

\reqfunccad
{REQ-CAD-TUR-01 Manter turma}
{O sistema deve manter o cadastro de turmas, que definem a recorrência semanal das aulas associando sala, horário de início, duração e professora responsável.}
{Essencial}
{Depende do requisito [\ref{REQ-CAD-SAL-01}].}
{\begin{itemize}[leftmargin=0.6cm]
    \item O sistema deve permitir cadastrar, consultar, atualizar, inativar e reativar turmas.
    \item O sistema deve exigir identificação, `horarioInicio', `duracao', `diaSemana' e `sala' como campos obrigatórios nos fluxos de inclusão e atualização.
    \item O sistema deve aceitar apenas valores válidos para `duracao' (positiva) e `diaSemana' (dentro da enumeração permitida).
    \item O sistema deve registrar a turma em estado inativo quando a ativação depender de validação posterior das regras de domínio.
    \item O sistema deve permitir consulta por sala, dia da semana, horário e status, com ordenação por dia da semana e horário de início.
    \item O sistema deve recusar ativação de turma sem sala válida, sem ao menos um dia de ocorrência ou com conflito de horário em relação a outra turma ativa na mesma sala.
\end{itemize}}

\rfcadastrotopico{Manter Manutenção}{REQ-CAD-MAN}{Essencial}

\reqfunccad
{REQ-CAD-MAN-01 Manter manutenção}
{O sistema deve manter o cadastro de manutenções, que registram intervenções em bikes e afetam sua disponibilidade operacional nas salas.}
{Essencial}
{Depende dos requisitos [\ref{REQ-CAD-BIK-01}] e [\ref{REQ-CAD-TMA-01}].}
{\begin{itemize}[leftmargin=0.6cm]
    \item O sistema deve permitir cadastrar, consultar, atualizar, cancelar, inativar e reativar registros de manutenção.
    \item O sistema deve exigir `bike', `tipoManutencao', `dataSolicitacao' e `descricao' como campos obrigatórios.
    \item O sistema deve permitir registrar manutenção sem `dataRealizacao'.
    \item O sistema deve gerenciar os estados da manutenção: pendente (sem `dataRealizacao' e sem cancelamento), realizada (com `dataRealizacao' válida igual ou posterior a `dataSolicitacao' e sem cancelamento), ou cancelada (independentemente de `dataRealizacao').
    \item O sistema deve gerar automaticamente código único e inalterável para cada novo registro de manutenção.
    \item O sistema deve permitir selecionar `bike' e `tipoManutencao' apenas dentre registros válidos e ativos apresentados pelo sistema.
    \item O sistema deve permitir consulta por `bike', `tipoManutencao', período de solicitação e estado operacional do registro.
    \item O sistema deve permitir filtros pelos estados \textbf{pendente}, \textbf{realizada} e \textbf{cancelada}.
    \item O sistema deve ordenar consultas por `dataSolicitacao' em ordem cronológica decrescente.
    \item O sistema deve refletir automaticamente no mapa da sala e na disponibilidade da turma os efeitos de manutenção pendente, realizada, cancelada ou associada a bike inativa.
    \item O sistema deve impedir que bikes com manutenção pendente apareçam como disponíveis.
\end{itemize}}

\rfmacrogrupo
{mg:gestao_aula}
{GDA}
{Gestão de aula}
{Professora}
{Essencial}
{Organizar a operação da aula, incluindo visualização e cancelamento administrativo de \textit{check-ins}, painéis de ocupação e frequência, e manter o repertório e os vínculos que definem seu uso por turma.}

\rfgruporeq
{g:operacao_aula}
{Operação da aula e painéis gerenciais}

\reqfuncfull
{rf:visualizar_mapa_operacional}
{Visualizar mapa operacional da turma}
{O sistema deve exibir à professora o mapa operacional da turma selecionada.}
{Importante}
{Depende dos requisitos [RF-AP-\ref{rf:restringir_acesso_perfil_professora}], [\ref{REQ-CAD-TUR-01}], [RF-CA-\ref{rf:realizar_checkin}] e [RF-CA-\ref{rf:fila_espera}].}
{\begin{itemize}[leftmargin=0.6cm]
    \item O sistema deve disponibilizar a representação da ocupação efetiva da turma na data consultada.
    \item O sistema deve associar os alunos às posições ocupadas, distinguindo os estados: confirmado, em fila de espera e cancelado.
    \item Quando não houver reservas ativas, o sistema deve exibir a configuração da sala sem nomes.
    \item Quando não houver sala válida ou contexto montável, o sistema deve informar a indisponibilidade.
    \item O sistema deve exibir discretamente a idade média dos alunos com \textit{check-in} ativo na turma; quando não houver nenhum \textit{check-in} ativo, o indicador não deve ser exibido.
\end{itemize}}

\reqfuncfull
{rf:cancelar_checkins_de_alunos}
{Cancelar \textit{check-ins} de alunos}
{O sistema deve cancelar check-ins de alunos da turma no contexto operacional, mediante ação da professora.}
{Importante}
{Depende dos requisitos [RF-AP-\ref{rf:restringir_acesso_perfil_professora}], [RF-GDA-\ref{rf:visualizar_mapa_operacional}] e [RF-CA-\ref{rf:cancelar_checkin}].}
{\begin{itemize}[leftmargin=0.6cm]
    \item Mediante ação da professora, o sistema deve permitir seleção apenas de posições ocupadas com \textit{check-in} ativo para cancelamento.
    \item O sistema deve impedir seleção de posições livres, em manutenção ou reservadas para a professora.
    \item Após o cancelamento, o sistema deve liberar a vaga imediatamente.
    \item O sistema deve manter o \textit{check-in} cancelado visível com o estado de cancelado.
    \item Se a reserva não estiver mais ativa no momento da confirmação, o sistema deve informar o estado atual e atualizar o mapa sem exibir sucesso indevido.
    \item O sistema deve recusar tentativas de cancelamento fora do contexto da turma ou do período permitido, exibindo mensagem compatível.
\end{itemize}}

\reqfuncfull
{rf:consultar_historico_de_mixes}
{Consultar histórico de mixes da turma}
{O sistema deve exibir à professora o histórico de \textit{mixes} utilizados pela turma.}
{Desejável}
{Depende dos requisitos [RF-AP-\ref{rf:restringir_acesso_perfil_professora}] e [\ref{REQ-CAD-TMX-01}].}
{\begin{itemize}[leftmargin=0.6cm]
    \item O sistema deve listar os \textit{mixes} encerrados da turma em ordem cronológica decrescente, exibindo nome e intervalo de uso de cada vínculo encerrado.
    \item Quando não houver histórico encerrado, o sistema deve informar a ausência e manter acessível o \textit{mix} vigente, quando houver.
    \item Quando a turma nunca tiver tido \textit{mix} associado, o sistema deve informar a ausência sem apresentar dados de outras turmas.
    \item O histórico deve refletir a vigência do vínculo `TurmaMix', não a vigência geral de `Mix'.
\end{itemize}}

\reqfuncfull
{rf:exibir_painel_ocupacao}
{Exibir painel de ocupação do dia}
{O sistema deve exibir à professora um painel de ocupação das turmas do dia atual.}
{Importante}
{Depende dos requisitos [RF-AP-\ref{rf:restringir_acesso_perfil_professora}], [\ref{REQ-CAD-TUR-01}], [RF-CA-\ref{rf:realizar_checkin}] e [RF-CA-\ref{rf:cancelar_checkin}].}
{\begin{itemize}[leftmargin=0.6cm]
    \item O sistema deve exibir todas as turmas ativas do dia com a respectiva taxa de ocupação, incluindo turmas sem \textit{check-ins} com ocupação zero.
    \item Quando não houver turmas ativas no dia, o sistema deve exibir mensagem informativa sem apresentar dados parciais.
    \item Ao selecionar uma turma, o sistema deve exibir o mapa operacional correspondente.
    \item Se a turma deixar de estar disponível entre a carga do painel e a seleção, o sistema deve informar a indisponibilidade e atualizar a listagem sem exibir dados desatualizados.
    \item Quando o sistema carregar o painel mas não conseguir calcular a ocupação de uma turma específica, o sistema deve exibir essa turma com indicação de dado indisponível e manter as demais turmas visíveis com seus dados calculados.
\end{itemize}}

\reqfuncfull
{rf:exibir_painel_frequencia}
{Exibir painel de frequência dos alunos}
{O sistema deve exibir à professora um painel de frequência individual dos alunos no período selecionado.}
{Importante}
{Depende dos requisitos [RF-AP-\ref{rf:restringir_acesso_perfil_professora}], [RF-CA-\ref{rf:realizar_checkin}] e [RF-CA-\ref{rf:cancelar_checkin}].}
{\begin{itemize}[leftmargin=0.6cm]
    \item O sistema deve suportar filtro por aluno, turma e período.
    \item A contagem de frequência deve incluir apenas \textit{check-ins} confirmados, excluindo cancelados.
    \item Cada linha do painel deve exibir o nome do aluno, o total de \textit{check-ins} confirmados e o percentual calculado exclusivamente sobre as aulas da turma filtrada no período selecionado.
    \item O sistema deve permitir ordenação dos resultados por nome ou por frequência.
    \item Quando não houver dados no período selecionado, o sistema deve exibir mensagem informativa sem apresentar linhas com valores zerados ou vazios.
    \item Quando o sistema carregar o painel mas não conseguir calcular a frequência para uma combinação específica de filtros, o sistema deve sinalizar o dado indisponível naquela linha sem ocultar os demais resultados calculados.
\end{itemize}}

\reqfuncfull
{rf:consultar_avaliacoes_musicas}
{Consultar avaliações das músicas}
{O sistema deve permitir que a professora consulte as avaliações registradas pelos alunos para as músicas do repertório.}
{Importante}
{Depende dos requisitos [RF-CA-\ref{rf:avaliar_musica_repertorio}], [\ref{REQ-CAD-MUS-01}] e [RF-AP-\ref{rf:restringir_acesso_perfil_professora}].}
{\begin{itemize}[leftmargin=0.6cm]
    \item O sistema deve permitir a consulta das avaliações associadas às músicas do repertório.
    \item O sistema deve apresentar, para cada música, a quantidade de avaliações registradas e a avaliação média correspondente.
    \item O sistema deve permitir consulta por música, artista ou banda e mix.
    \item O sistema deve permitir ordenar os resultados pela avaliação média ou pela quantidade de avaliações registradas.
    \item Quando uma música não possuir avaliações, o sistema deve apresentar essa informação sem atribuir nota média artificial.
    \item Quando não houver resultados compatíveis com os filtros aplicados, o sistema deve exibir mensagem informativa.
\end{itemize}}

\reqfuncfull
{rf:copiar_marcacoes_instagram_alunos_turma}
{Copiar marcações de alunos para publicação no Instagram}
{O sistema deve permitir que a professora copie a relação de perfis do Instagram dos alunos vinculados a uma turma, para utilização em publicações realizadas fora do sistema.}
{Desejável}
{Depende dos requisitos [\ref{REQ-CAD-TUR-01}], [\ref{REQ-CAD-ALU-01}] e [RF-AP-\ref{rf:restringir_acesso_perfil_professora}].}
{\begin{itemize}[leftmargin=0.6cm]
    \item O sistema deve disponibilizar a funcionalidade a partir da consulta de uma turma.
    \item O sistema deve apresentar uma opção específica para visualizar a relação de perfis do Instagram dos alunos vinculados à turma selecionada.
    \item O sistema deve incluir somente alunos que possuam perfil do Instagram informado em seu cadastro.
    \item O sistema deve apresentar os perfis em formato textual adequado para marcação em publicações, utilizando o prefixo \texttt{@}.
    \item O sistema deve permitir copiar integralmente a relação de perfis para a área de transferência do dispositivo.
    \item O sistema deve informar à professora quando a cópia for concluída com sucesso.
    \item Quando nenhum aluno da turma possuir perfil do Instagram informado, o sistema deve apresentar mensagem informativa.
\end{itemize}}

\rfgruporeq
{g:cadastros_repertorio}
{Cadastros musicais do repertório}

\rfcadastrotopico{Manter Artista ou Banda}{REQ-CAD-ARB}{Importante}

\reqfunccad
{REQ-CAD-ARB-01 Manter artista ou banda}
{O sistema deve manter o cadastro de artistas e bandas, referenciados nas músicas do repertório como responsáveis pelas obras.}
{Importante}
{Nenhuma.}
{\begin{itemize}[leftmargin=0.6cm]
    \item O sistema deve permitir cadastrar, consultar, atualizar, inativar e reativar artistas e bandas.
    \item O sistema deve exigir `nome' como campo obrigatório nos fluxos de inclusão e atualização.
    \item O sistema deve permitir cadastro sem `descricao', `link' ou `urlFoto'.
    \item O sistema deve validar formato de `link' e `urlFoto' quando informados, e unicidade de `nome' na inclusão e atualização.
    \item O sistema deve permitir consulta por `nome' e status, com ordenação alfabética por nome.
\end{itemize}}

\rfcadastrotopico{Manter Categoria Musical}{REQ-CAD-CAT}{Importante}

\reqfunccad
{REQ-CAD-CAT-01 Manter categoria musical}
{O sistema deve manter categorias musicais e suas associações às músicas cadastradas, permitindo classificar o repertório por estilo ou gênero sem sobrecarregar o cadastro principal de música.}
{Importante}
{Depende do requisito [\ref{REQ-CAD-MUS-01}].}
{\begin{itemize}[leftmargin=0.6cm]
    \item O sistema deve permitir cadastrar, consultar, atualizar, inativar e reativar categorias musicais.
    \item O sistema deve permitir buscar uma música cadastrada por digitação, apresentando opções correspondentes para seleção controlada.
    \item O sistema deve exigir `nome' como campo obrigatório nos fluxos de inclusão e atualização.
    \item O sistema deve permitir cadastro sem `descricao'.
    \item O sistema deve validar unicidade de `nome' na inclusão e na atualização.
    \item O sistema deve permitir selecionar múltiplas categorias para a música escolhida, adicionando cada categoria selecionada a um painel de resultados da seleção.
    \item Quando a categoria digitada não existir, o sistema deve permitir cadastrá-la como nova categoria antes de associá-la à música.
    \item O sistema deve permitir remover categorias do painel de seleção antes da confirmação.
    \item O sistema não deve permitir associar a mesma categoria musical mais de uma vez à mesma música.
    \item O sistema deve permitir consulta por `nome' e status, com ordenação alfabética por nome.
\end{itemize}}

\rfcadastrotopico{Manter Vídeo de Aula}{REQ-CAD-VID}{Importante}

\reqfunccad
{REQ-CAD-VID-01 Manter vídeo de aula}
{O sistema deve manter links de vídeos de aula e suas associações às músicas cadastradas, para consulta pelos alunos durante as aulas.}
{Importante}
{Depende do requisito [\ref{REQ-CAD-MUS-01}].}
{\begin{itemize}[leftmargin=0.6cm]
    \item O sistema deve permitir cadastrar, consultar, atualizar, inativar e reativar vídeos de aula.
    \item O sistema deve permitir buscar uma música cadastrada por digitação, apresentando opções correspondentes para seleção controlada.
    \item O sistema deve exigir `nome' e `linkVideo' como campos obrigatórios nos fluxos de inclusão e atualização, sendo `linkVideo' uma URL para vídeo externo, não envio de arquivo.
    \item O sistema deve validar formato de `linkVideo'.
    \item O sistema deve validar unicidade de `linkVideo' na inclusão e na atualização.
    \item O sistema deve permitir selecionar múltiplos links de vídeo para a música escolhida, adicionando cada link selecionado a um painel de resultados da seleção.
    \item Quando a URL digitada não existir, o sistema deve permitir cadastrá-la como novo vídeo de aula antes de associá-la à música.
    \item O sistema deve permitir remover vídeos de aula do painel de seleção antes da confirmação.
    \item O sistema não deve permitir associar o mesmo vídeo de aula mais de uma vez à mesma música.
    \item O sistema deve permitir consulta por `nome', `linkVideo' e status, com ordenação alfabética por nome.
\end{itemize}}

\rfcadastrotopico{Manter Música}{REQ-CAD-MUS}{Essencial}

\reqfunccad
{REQ-CAD-MUS-01 Manter música}
{O sistema deve manter o cadastro de músicas, que compõem o repertório dos mixes e são exibidas aos alunos durante a consulta da aula.}
{Essencial}
{Depende do requisito [\ref{REQ-CAD-ARB-01}].}
{\begin{itemize}[leftmargin=0.6cm]
    \item O sistema deve permitir cadastrar, consultar, atualizar, inativar e reativar músicas.
    \item O sistema deve exigir `nome' e `artistaBanda' como campos obrigatórios na inclusão e na atualização.
    \item O sistema deve permitir cadastrar música sem categorias musicais e sem vídeos de aula associados, pois essas associações são mantidas em fluxos próprios.
    \item O sistema deve validar a unicidade da combinação de `nome' e `artistaBanda' na inclusão e na atualização.
    \item O sistema deve permitir selecionar `artistaBanda' por seleção única controlada.
    \item O sistema deve permitir consulta por `nome', `artistaBanda', `categoriaMusical' e status, com ordenação alfabética por nome.
    \item O sistema deve manter mixes já cadastrados exibindo a referência da música após atualizações válidas ou inativação.
\end{itemize}}

\rfcadastrotopico{Manter Mix}{REQ-CAD-MIX}{Essencial}

\reqfunccad
{REQ-CAD-MIX-01 Manter mix}
{O sistema deve manter o cadastro de mixes, que são coleções ordenadas de músicas vinculadas a turmas por períodos de vigência para uso nas aulas.}
{Essencial}
{Depende do requisito [\ref{REQ-CAD-MUS-01}].}
{\begin{itemize}[leftmargin=0.6cm]
    \item O sistema deve permitir cadastrar, consultar, atualizar, inativar e reativar mixes.
    \item O sistema deve exigir `nome', `descricao' e ao menos uma música com ordem definida como campos obrigatórios.
    \item O sistema deve exigir ordem completa e inequívoca entre as músicas associadas na inclusão e na atualização.
    \item O sistema deve permitir cadastrar mix sem `dataFim'.
    \item O sistema deve validar unicidade de `nome', e que `dataFim' seja igual ou posterior a `dataInicio' quando informada, na inclusão e atualização.
    \item O sistema deve permitir selecionar músicas e definir a ordem de execução por controles da interface.
    \item O sistema deve permitir consulta por `nome', período de vigência e status, com ordenação cronológica decrescente por `dataInicio'.
\end{itemize}}

\rfgruporeq
{g:vinculos_vigencia_aula}
{Vínculo entre turma e mix e vigência do repertório}

\rfcadastrotopico{Manter Vínculo TurmaMix}{REQ-CAD-TMX}{Essencial}

\reqfunccad
{REQ-CAD-TMX-01 Manter vínculo TurmaMix}
{O sistema deve manter os vínculos entre turma e mix, que determinam qual repertório está ativo em cada turma em cada período de vigência.}
{Essencial}
{Depende dos requisitos [\ref{REQ-CAD-TUR-01}] e [\ref{REQ-CAD-MIX-01}].}
{\begin{itemize}[leftmargin=0.6cm]
    \item O sistema deve permitir cadastrar, consultar, atualizar, encerrar, inativar e reativar vínculos TurmaMix.
    \item O sistema deve exigir `turma', `mix' e `dataInicio' como campos obrigatórios.
    \item O sistema deve permitir cadastrar vínculo sem `dataFim'.
    \item O sistema deve exigir que `dataFim' seja igual ou posterior a `dataInicio' quando informada.
    \item O sistema não deve permitir que a mesma turma possua mais de um vínculo ativo de mix simultaneamente.
    \item O sistema deve permitir consulta por turma, mix, período e status, com ordenação cronológica decrescente por `dataInicio'.
    \item O sistema deve recuperar vínculos ativos, inativos e encerrados nas consultas.
    \item O sistema deve permitir identificar com clareza o mix vigente da turma e seus períodos anteriores, sem ambiguidade.
\end{itemize}}

\rfmacrogrupo
{mg:acompanhamento_aluno}
{AA}
{Acompanhamento do Aluno}
{Aluno}
{Importante}
{Nesta versão, o acompanhamento individual do aluno cobre a consulta ao próprio histórico de \textit{check-ins}, com acesso restrito ao aluno autenticado.}

\rfgruporeqatorprio
{g:historico_metricas_aluno}
{Histórico individual de participação}
{Aluno}
{Desejável}

\reqfuncprio
{rf:consultar_historico_de_presenca}
{Consultar histórico de presença}
{O sistema deve exibir ao aluno autenticado o próprio histórico de \textit{check-ins}.}
{Importante}
{Depende dos requisitos [RF-CA-\ref{rf:realizar_checkin}] e [RF-CA-\ref{rf:cancelar_checkin}].}
{\begin{itemize}[leftmargin=0.6cm]
    \item O sistema deve reunir apenas os \textit{check-ins} do aluno autenticado na consulta.
    \item O sistema deve distinguir \textit{check-ins} futuros de concluídos e exibir cancelados com o respectivo estado.
    \item Quando não houver histórico, o sistema deve exibir mensagem informativa.
    \item O sistema deve recusar acesso ao histórico sem autenticação válida.
    \item O sistema não deve exibir o histórico de \textit{check-ins} de terceiros.
\end{itemize}}

\subsection{Matriz de Dependências dos Requisitos Funcionais}

A Tabela~\ref{tab:matriz_dependencia} consolida as dependências declaradas entre os requisitos funcionais, após a verificação de consistência entre os cadastros, os fluxos de interação do aluno, as operações realizadas pela professora e as funcionalidades de acompanhamento individual.

Cada requisito também foi registrado como uma \textit{Issue} no GitHub. O acesso à respectiva \textit{Issue} pode ser realizado por meio do identificador apresentado na coluna \textit{Issue} da tabela, substituindo-se \texttt{XX} pelo número correspondente no seguinte endereço:

\url{https://github.com/heliokamakawa/spin_flow_pt_br/issues/XX}

\begingroup
\footnotesize
\begin{longtable}{|c|p{6cm}|p{7cm}|} 
    \caption{Matriz de Dependências dos Requisitos Funcionais.}
    \label{tab:matriz_dependencia} \\ \hline
    \textbf{Issue} & \textbf{Requisito} & \textbf{Dependência de requisitos} \\ \hline
    \endfirsthead 
    
    \caption[]{Matriz de Dependências dos Requisitos Funcionais (continuação).} \\ \hline
    \textbf{Issue} & \textbf{Requisito} & \textbf{Dependência de requisitos} \\ \hline
    \endhead

    1 & [RF-AP-\ref{rf:autenticar_login_aluno}] Autenticar login do aluno & [\ref{REQ-CAD-ALU-01}]. \\ \hline
    2 & [RF-AP-\ref{rf:recuperar_senha_aluno}] Recuperar senha do aluno & [\ref{REQ-CAD-ALU-01}]. \\ \hline
    3 & [RF-AP-\ref{rf:recuperar_acesso_aluno_dispositivo}] Recuperar acesso do aluno em novo dispositivo & [\ref{REQ-CAD-ALU-01}]. \\ \hline
    4 & [RF-AP-\ref{rf:restringir_acesso_perfil_aluno}] Restringir acesso do perfil de aluno & [RF-AP-\ref{rf:autenticar_login_aluno}]. \\ \hline
    5 & [RF-AP-\ref{rf:direcionar_login_aluno}] Direcionar login do aluno & [RF-AP-\ref{rf:autenticar_login_aluno}]. \\ \hline
    6 & [RF-AP-\ref{rf:autenticar_login_professora}] Autenticar login da professora & Nenhuma. \\ \hline
    7 & [RF-AP-\ref{rf:restringir_acesso_perfil_professora}] Restringir acesso do perfil de professora & [RF-AP-\ref{rf:autenticar_login_professora}]. \\ \hline
    8 & [RF-AP-\ref{rf:direcionar_login_professora}] Direcionar login da professora & [RF-AP-\ref{rf:autenticar_login_professora}]. \\ \hline
    9 & [RF-CA-\ref{rf:painel_aulas_dia}] Painel geral de aulas do dia & [\ref{REQ-CAD-TUR-01}] e [\ref{REQ-CAD-SAL-01}]. \\ \hline
    10 & [RF-CA-\ref{rf:detalhe_aula_selecionada}] Exibir detalhes da aula selecionada & [RF-CA-\ref{rf:painel_aulas_dia}], [\ref{REQ-CAD-BIK-01}] e [\ref{REQ-CAD-MAN-01}]. \\ \hline
    11 & [RF-CA-\ref{rf:realizar_checkin}] Realizar \textit{check-in} & [\ref{REQ-CAD-ALU-01}] e [RF-CA-\ref{rf:detalhe_aula_selecionada}]. \\ \hline
    12 & [RF-CA-\ref{rf:cancelar_checkin}] Cancelar \textit{check-in} & [RF-CA-\ref{rf:realizar_checkin}]. \\ \hline
    13 & [RF-CA-\ref{rf:fila_espera}] Entrar na fila de espera & [\ref{REQ-CAD-ALU-01}], [RF-CA-\ref{rf:detalhe_aula_selecionada}] e [RF-CA-\ref{rf:realizar_checkin}]. \\ \hline
    14 & [RF-CA-\ref{rf:consultar_repertorio_aula}] Consultar repertório e detalhes musicais da aula & [\ref{REQ-CAD-MIX-01}], \ref{REQ-CAD-TMX-01}] e [RF-CA-\ref{rf:detalhe_aula_selecionada}]. \\ \hline
    15 & [RF-CA-\ref{rf:acesso_painel_participacao}] Acesso ao painel de participação do aluno & [RF-AP-\ref{rf:autenticar_login_aluno}] e [RF-AA-\ref{rf:consultar_historico_de_presenca}]. \\ \hline
    44 & [RF-CA-\ref{rf:avaliar_musica_repertorio}] Registrar e modificar avaliação de música & [RF-CA-\ref{rf:consultar_repertorio_aula}] e [\ref{REQ-CAD-MUS-01}]. \\ \hline
    16 & [\ref{REQ-CAD-ALU-01}] Manter aluno & Nenhuma. \\ \hline
    17 & [\ref{REQ-CAD-GAL-01}] Manter grupo de alunos & [\ref{REQ-CAD-ALU-01}]. \\ \hline
    18 & [\ref{REQ-CAD-FAB-01}] Manter fabricante & Nenhuma. \\ \hline
    19 & [\ref{REQ-CAD-TMA-01}] Manter tipo de manutenção & Nenhuma. \\ \hline
    20 & [\ref{REQ-CAD-SAL-01}] Manter sala & Nenhuma. \\ \hline
    21 & [\ref{REQ-CAD-BIK-01}] Manter \textit{bike} & [\ref{REQ-CAD-FAB-01}] e [\ref{REQ-CAD-SAL-01}]. \\ \hline
    22 & [\ref{REQ-CAD-TUR-01}] Manter turma & [\ref{REQ-CAD-SAL-01}]. \\ \hline
    23 & [\ref{REQ-CAD-MAN-01}] Manter manutenção & [\ref{REQ-CAD-BIK-01}] e [\ref{REQ-CAD-TMA-01}]. \\ \hline
    24 & [RF-GDA-\ref{rf:visualizar_mapa_operacional}] Visualizar mapa operacional da turma & [RF-AP-\ref{rf:restringir_acesso_perfil_professora}], [\ref{REQ-CAD-TUR-01}], [RF-CA-\ref{rf:realizar_checkin}] e [RF-CA-\ref{rf:fila_espera}]. \\ \hline
    25 & [RF-GDA-\ref{rf:cancelar_checkins_de_alunos}] Cancelar \textit{check-ins} de alunos & [RF-AP-\ref{rf:restringir_acesso_perfil_professora}], [RF-GDA-\ref{rf:visualizar_mapa_operacional}] e [RF-CA-\ref{rf:cancelar_checkin}]. \\ \hline
    26 & [RF-GDA-\ref{rf:consultar_historico_de_mixes}] Consultar histórico de mixes da turma & [RF-AP-\ref{rf:restringir_acesso_perfil_professora}] e [\ref{REQ-CAD-TMX-01}]. \\ \hline
    27 & [RF-GDA-\ref{rf:exibir_painel_ocupacao}] Exibir painel de ocupação do dia & [RF-AP-\ref{rf:restringir_acesso_perfil_professora}], [\ref{REQ-CAD-TUR-01}], [RF-CA-\ref{rf:realizar_checkin}] e [RF-CA-\ref{rf:cancelar_checkin}]. \\ \hline
    28 & [RF-GDA-\ref{rf:exibir_painel_frequencia}] Exibir painel de frequência dos alunos & [RF-AP-\ref{rf:restringir_acesso_perfil_professora}], [RF-CA-\ref{rf:realizar_checkin}] e [RF-CA-\ref{rf:cancelar_checkin}]. \\ \hline
    45 & [RF-GDA-\ref{rf:consultar_avaliacoes_musicas}] Consultar avaliações das músicas & [RF-CA-\ref{rf:avaliar_musica_repertorio}], [\ref{REQ-CAD-MUS-01}] e [RF-AP-\ref{rf:restringir_acesso_perfil_professora}]. \\ \hline
    46 & [RF-GDA-\ref{rf:copiar_marcacoes_instagram_alunos_turma}] Copiar marcações de alunos para publicação no Instagram &  [\ref{REQ-CAD-TUR-01}], [\ref{REQ-CAD-ALU-01}] e [RF-AP-\ref{rf:restringir_acesso_perfil_professora}]. \\ \hline
    29 & [\ref{REQ-CAD-ARB-01}] Manter artista ou banda & Nenhuma. \\ \hline
    30 & [\ref{REQ-CAD-CAT-01}] Manter categoria musical & [\ref{REQ-CAD-MUS-01}]. \\ \hline
    31 & [\ref{REQ-CAD-VID-01}] Manter vídeo de aula & [\ref{REQ-CAD-MUS-01}]. \\ \hline
    32 & [\ref{REQ-CAD-MUS-01}] Manter música & [\ref{REQ-CAD-ARB-01}]. \\ \hline
    33 & [\ref{REQ-CAD-MIX-01}] Manter mix & [\ref{REQ-CAD-MUS-01}]. \\ \hline
    34 & [\ref{REQ-CAD-TMX-01}] Manter vínculo TurmaMix & [\ref{REQ-CAD-TUR-01}] e [\ref{REQ-CAD-MIX-01}]. \\ \hline
    35 & [RF-AA-\ref{rf:consultar_historico_de_presenca}] Consultar histórico de presença & [RF-CA-\ref{rf:realizar_checkin}] e [RF-CA-\ref{rf:cancelar_checkin}]. \\ \hline
\end{longtable}
\endgroup